% !TeX root = part6.tex
\documentclass[../final.tex]{subfiles}

\begin{document}

% Лекция 06.10.2026.
% Источник: Desktop/1.jpg (с отметки 06.10), 2.jpg, 3.jpg,
% 4.jpg (до отметки 13.10). Продолжение главы об ОДУ.
\ifSubfilesClassLoaded{\setcounter{chapter}{5}\setcounter{section}{3}}{}
\captionsetup{font=footnotesize,labelfont=bf,skip=5pt}

\subsection{Невязка модифицированного метода Эйлера}

В предыдущей лекции значение решения в середине шага заменялось
полушагом Эйлера. Полученная схема
\begin{equation*}
    y_{i+1}=y_i+h f\!\left(x_i+\frac h2,
    y_i+\frac h2 f(x_i,y_i)\right)
\end{equation*}
называется также \rconcept{явным методом средней точки}.
Проверим её порядок через невязку, то есть результат подстановки
точного решения $u_i=u(x_i)$ в разностное уравнение:
\begin{equation*}
    \Psi_{h,i}=\frac{u_{i+1}-u_i}{h}
    -f\!\left(x_i+\frac h2,u_i+\frac h2 f(x_i,u_i)\right).
\end{equation*}
Далее для краткости $f$, $f_x$ и $f_u$ без аргументов означают значения
в точке $(x_i,u_i)$. По формуле Тейлора
\begin{align*}
    \frac{u_{i+1}-u_i}{h}
        &=u'_i+\frac h2u''_i+O(h^2),\\
    f\!\left(x_i+\frac h2,u_i+\frac h2f\right)
        &=f+\frac h2f_x+\frac h2f f_u+O(h^2).
\end{align*}
Так как $u'_i=f$ и $u''_i=f_x+f f_u$, члены нулевого и первого
порядков сокращаются:
\begin{equation*}
    \Psi_{h,i}
    =u'_i-f+\frac h2\bigl(u''_i-f_x-f f_u\bigr)+O(h^2)
    =O(h^2).
\end{equation*}
Следовательно, схема имеет \rresult{второй порядок аппроксимации}.
Погрешность одного шага, начатого с точного значения, равна
$h\Psi_{h,i}=O(h^3)$. Глобальная погрешность на фиксированном отрезке
при условиях гладкости и Липшица имеет порядок $O(h^2)$.

\section{Исправленный метод Эйлера}

\subsection{От формулы трапеций к схеме предиктор--корректор}

Запишем точное приращение решения в интегральной форме:
\begin{equation*}
    u(x_{i+1})-u(x_i)
    =\int_{x_i}^{x_{i+1}}f(x,u(x))\,dx.
\end{equation*}
Заменяя интеграл формулой трапеций, получаем
\begin{equation*}
    u(x_{i+1})=u(x_i)+\frac h2
    \bigl[f(x_i,u(x_i))+f(x_{i+1},u(x_{i+1}))\bigr]+O(h^3).
\end{equation*}
Если перенести эту формулу непосредственно на сеточные значения, получится
неявная схема: $y_{i+1}$ входит в правую часть. Чтобы получить явный
алгоритм, сначала предскажем значение на конце шага методом Эйлера,
а затем используем этот прогноз для вычисления второго наклона:
\begin{equation*}
    \begin{cases}
        \widetilde y_{i+1}=y_i+h f(x_i,y_i),\\[0.4em]
        y_{i+1}=y_i+\dfrac h2
        \bigl[f(x_i,y_i)+f(x_i+h,\widetilde y_{i+1})\bigr].
    \end{cases}
\end{equation*}
Первая строка --- \rconcept{предиктор}, вторая ---
\rconcept{корректор}. В одной формуле исправленный метод Эйлера
записывается как
\begin{equation*}
    y_{i+1}=y_i+\frac h2\left[
    f(x_i,y_i)+f\bigl(x_i+h,y_i+h f(x_i,y_i)\bigr)\right].
\end{equation*}
Геометрически берётся среднее двух наклонов: в начале шага и
в предсказанной конечной точке.

\subsection{Проверка порядка}

При $y_i=u_i$ разложение второго наклона даёт
\begin{equation*}
    f(x_i+h,u_i+hf)=f+h(f_x+f f_u)+O(h^2).
\end{equation*}
Поэтому
\begin{align*}
    y_{i+1}
    &=u_i+\frac h2\bigl[2f+h(f_x+f f_u)+O(h^2)\bigr]\\
    &=u_i+hu'_i+\frac{h^2}{2}u''_i+O(h^3).
\end{align*}
Разложение совпадает с точным решением до членов порядка $h^2$.
Исправленный метод Эйлера, как и метод средней точки, имеет
\rresult{второй порядок точности}, но вычисляет второй наклон
в другой точке.

\section{Методы Рунге--Кутты второго порядка}

\subsection{Общая двухстадийная схема}

Оба варианта Эйлера можно включить в одно семейство:
\begin{equation*}
    \begin{cases}
        k_1=f(x_i,y_i),\\
        k_2=f(x_i+b_1h,y_i+b_2h k_1),\\
        y_{i+1}=y_i+h(a_1 k_1+a_2 k_2).
    \end{cases}
\end{equation*}
Величины $k_1$ и $k_2$ --- \rconcept{стадии метода}, то есть
два значения правой части. Коэффициенты $a_1,a_2,b_1,b_2$
подбираются так, чтобы невязка была $O(h^2)$.

Подставим точное решение и разложим вторую стадию:
\begin{equation*}
    f(x_i+b_1h,u_i+b_2hf)
    =f+h(b_1f_x+b_2f f_u)+O(h^2).
\end{equation*}
Тогда невязка равна
\begin{align*}
    \Psi_{h,i}
    &=\frac{u_{i+1}-u_i}{h}
        -a_1f-a_2 f(x_i+b_1h,u_i+b_2hf)\\
    &=(1-a_1-a_2)f
      +h\left(\frac12-a_2b_1\right)f_x
      +h\left(\frac12-a_2b_2\right)f f_u+O(h^2).
\end{align*}
Чтобы сокращение происходило для произвольной гладкой правой части,
коэффициенты перед независимыми выражениями должны обращаться в нуль:
\begin{equation*}
    a_1+a_2=1,\qquad a_2b_1=\frac12,\qquad a_2b_2=\frac12.
\end{equation*}
Это \rconcept{условия второго порядка}.

\subsection{Однопараметрическое семейство}

Положим $a_2=\alpha\ne0$. Тогда
\begin{equation*}
    a_1=1-\alpha,\qquad b_1=b_2=\frac{1}{2\alpha},
\end{equation*}
и схема принимает вид
\begin{equation*}
    y_{i+1}=y_i+h\left[(1-\alpha)f(x_i,y_i)
    +\alpha f\!\left(x_i+\frac{h}{2\alpha},
    y_i+\frac{h}{2\alpha}f(x_i,y_i)\right)\right].
\end{equation*}
Таким образом, второй порядок не определяет метод единственным образом.
В частности,
\begin{itemize}
    \item $\alpha=1$: $a_1=0$, $b_1=b_2=1/2$ ---
    модифицированный метод Эйлера, или явный метод средней точки;
    \item $\alpha=1/2$: $a_1=a_2=1/2$, $b_1=b_2=1$ ---
    исправленный метод Эйлера;
    \item $\alpha=3/4$: $a_1=1/4$, $b_1=b_2=2/3$ ---
    метод Ралстона.
\end{itemize}
Последний вариант имеет вид
\begin{equation*}
    y_{i+1}=y_i+\frac h4\left[
    f(x_i,y_i)+3f\!\left(x_i+\frac{2h}{3},
    y_i+\frac{2h}{3}f(x_i,y_i)\right)\right].
\end{equation*}

\remark
В лекции выбор $\alpha=3/4$ отмечен как вариант с минимизацией
ошибки. Здесь речь идёт о выборе коэффициентов по оценке ведущего
члена локальной погрешности. \rnote{Универсального минимума фактической
ошибки для любой функции $f$ это не означает}: все перечисленные схемы
имеют второй порядок, а коэффициент ошибки зависит от уравнения.

\section{Классический метод Рунге--Кутты четвёртого порядка}

\subsection{Четыре вычисления правой части}

Формула Симпсона использует значения подынтегральной функции в начале,
середине и конце интервала:
\begin{equation*}
    \int_{x_i}^{x_i+h}g(x)\,dx
    \simeq\frac h6\left[g(x_i)+4g\!\left(x_i+\frac h2\right)
    +g(x_i+h)\right].
\end{equation*}
При решении ОДУ функция $g(x)=f(x,u(x))$ зависит от неизвестного
решения. Поэтому значения на промежуточных точках приходится
строить последовательно. \rconcept{Классическая схема РК4}:
\begin{equation*}
    \begin{cases}
        k_1=f(x_i,y_i),\\[0.35em]
        k_2=f\!\left(x_i+\dfrac h2,y_i+\dfrac h2 k_1\right),\\[0.35em]
        k_3=f\!\left(x_i+\dfrac h2,y_i+\dfrac h2 k_2\right),\\[0.35em]
        k_4=f(x_i+h,y_i+h k_3),\\[0.35em]
        y_{i+1}=y_i+\dfrac h6(k_1+2k_2+2k_3+k_4).
    \end{cases}
\end{equation*}
Здесь $k_j$ имеют смысл производных; множитель $h$ в них не включён.
На один шаг требуются четыре вычисления $f$.

Вторая стадия оценивает наклон в середине шага по прогнозу из $k_1$,
третья уточняет этот наклон по $k_2$, а четвёртая оценивает наклон
на конце шага по $k_3$. После этого наклоны складываются с весами
$1,2,2,1$.

\subsection{Связь с квадратурой и порядок точности}

Если правая часть зависит только от $x$, то $k_2=k_3=f(x_i+h/2)$,
и РК4 действительно превращается в формулу Симпсона.
При зависимости от $u$ две средние стадии различны:
\rnote{простая подстановка одного прогноза в формулу Симпсона
сама по себе не доказывает четвёртый порядок}.

Для достаточно гладкой $f$ согласование разложения схемы с рядом Тейлора
точного решения до $h^4$ даёт локальную погрешность $O(h^5)$.
На фиксированном отрезке при условии Липшица глобальная погрешность
равна $O(h^4)$. Поэтому РК4 имеет \rresult{четвёртый порядок точности}.

\section{Многошаговый метод Адамса--Башфорта}

\subsection{Интерполяция правой части по предыдущим шагам}

В одношаговых методах новое значение строится по текущему $y_i$,
хотя внутри шага может вычисляться несколько стадий. В
\rconcept{многошаговых методах} используются также предыдущие
сеточные значения.

Пусть уже известны $y_i$, $y_{i-1}$ и $y_{i-2}$ на равномерной сетке.
Обозначим
\begin{equation*}
    F_j=f(x_j,y_j).
\end{equation*}
Построим квадратичный полином по трём последним значениям правой части.
В форме Ньютона он записывается как
\begin{equation*}
    P_2(x)=F_i+F[x_i,x_{i-1}](x-x_i)
    +F[x_i,x_{i-1},x_{i-2}](x-x_i)(x-x_{i-1}),
\end{equation*}
где квадратные скобки обозначают разделённые разности. На равномерной сетке
\begin{align*}
    F[x_i,x_{i-1}]&=\frac{F_i-F_{i-1}}h,\\
    F[x_i,x_{i-1},x_{i-2}]
        &=\frac{F_i-2F_{i-1}+F_{i-2}}{2h^2}.
\end{align*}
Полином строится по прошлым точкам, а интегрируется на следующем
интервале $[x_i,x_{i+1}]$. Это экстраполяция правой части на один шаг.

\subsection{Интегрирование полинома}

Заменим $f(x,u(x))$ в интегральной формуле полиномом $P_2(x)$:
\begin{equation*}
    y_{i+1}=y_i+\int_{x_i}^{x_i+h}P_2(x)\,dx.
\end{equation*}
Положим $s=x-x_i$. Тогда $x-x_{i-1}=s+h$, и
\begin{align*}
    \int_0^h ds&=h,\qquad \int_0^h s\,ds=\frac{h^2}{2},\\
    \int_0^h s(s+h)\,ds
        &=\frac{h^3}{3}+\frac{h^3}{2}=\frac{5h^3}{6}.
\end{align*}
Отсюда
\begin{align*}
    y_{i+1}
    &=y_i+hF_i+\frac{h^2}{2}F[x_i,x_{i-1}]
        +\frac{5h^3}{6}F[x_i,x_{i-1},x_{i-2}]\\
    &=y_i+h\left[F_i+\frac12(F_i-F_{i-1})
        +\frac5{12}(F_i-2F_{i-1}+F_{i-2})\right].
\end{align*}
Собирая коэффициенты, получаем
\begin{equation*}
    y_{i+1}=y_i+\frac h{12}
    \left(23F_i-16F_{i-1}+5F_{i-2}\right),\qquad i\ge2.
\end{equation*}
Это \rconcept{трёхшаговый метод Адамса--Башфорта третьего порядка}.

\subsection{Погрешность и начало расчёта}

Для точной функции $F(x)=f(x,u(x))$ остаток квадратичной интерполяции
на интервале следующего шага имеет порядок $O(h^3)$:
\begin{equation*}
    F(x)-P_2(x)
    =\frac{F^{(3)}(\xi)}{3!}
    (x-x_i)(x-x_{i-1})(x-x_{i-2}).
\end{equation*}
После интегрирования на отрезке длины $h$ получаем локальную
погрешность $O(h^4)$. При гладкости правой части и достаточно точных
начальных значениях глобальная погрешность составляет $O(h^3)$.

Для старта недостаточно одного условия $y_0=u_0$: нужны ещё $y_1$ и $y_2$.
Их можно найти, например, двумя шагами РК4. После запуска достаточно
одного нового вычисления $F_i$ на шаг: два предыдущих значения
сохраняются из прошлых шагов.

\clearpage
\section{Устойчивость решения и численной схемы}

\subsection{Устойчивость по Ляпунову}

Рассмотрим задачу Коши
\begin{equation*}
    u'=f(x,u),\qquad u(a)=u_0.
\end{equation*}
Пусть $u(x)$ --- её решение, а $v(x)$ --- решение того же уравнения
с немного изменённым начальным значением.
Решение $u$ называется \rconcept{устойчивым по Ляпунову}, если
для любого $\varepsilon>0$ существует $\delta>0$ такое, что
\begin{equation*}
    \|v(a)-u(a)\|<\delta
    \quad\Longrightarrow\quad
    \|v(x)-u(x)\|<\varepsilon\quad\text{для всех }x\ge a.
\end{equation*}
Иными словами, достаточно малое начальное возмущение остаётся малым
на всём последующем промежутке. Здесь $\|\cdot\|$ обозначает норму;
для одного уравнения это обычный модуль.

Если, кроме устойчивости, для достаточно близких начальных данных
выполняется
\begin{equation*}
    \|v(x)-u(x)\|\longrightarrow0,\qquad x\longrightarrow\infty,
\end{equation*}
то решение называется \rconcept{асимптотически устойчивым}.

\subsection{Линейная система и собственные значения}

Для системы с постоянной матрицей
\begin{equation*}
    \frac{d\boldsymbol u}{dx}=A\boldsymbol u,\qquad
    \boldsymbol u(a)=\boldsymbol u_0
\end{equation*}
возмущение $\boldsymbol\delta=\boldsymbol v-\boldsymbol u$
подчиняется тому же уравнению:
\begin{equation*}
    \boldsymbol\delta'=A\boldsymbol\delta,\qquad
    \boldsymbol\delta(x)=e^{A(x-a)}\boldsymbol\delta(a).
\end{equation*}
Если $\boldsymbol q_j$ --- собственный вектор $A$ с собственным
значением $\mu_j$, то соответствующая мода изменяется как
$e^{\mu_j(x-a)}$. Поэтому условие
\begin{equation*}
    \operatorname{Re}\mu_j<0\quad\text{для всех }j
\end{equation*}
обеспечивает \rresult{асимптотическую устойчивость линейной системы}.
Если хотя бы одно собственное значение имеет положительную
вещественную часть, соответствующее возмущение растёт.

\remark
Случай $\operatorname{Re}\mu_j=0$ требует отдельного анализа:
жордановы блоки могут давать степенной рост даже без положительной
вещественной части собственных значений.

Численная схема также должна корректно переносить малые возмущения
данных. В отличие от непрерывного уравнения, их распространение
зависит ещё и от шага $h$ и от выбранного метода. В следующей лекции
это исследуется на затухающем линейном уравнении.

\end{document}
